\notepage{12}{Degenerate eigenspaces and diagonalization}
\smallheading{Consequences of Hermiticity}
Setting $i=j$ in
$(\omega_i-\omega_j^*)\braket{\omega_j}{\omega_i}=0$ and using
$\braket{\omega_i}{\omega_i}>0$ shows that
\[
 \omega_i=\omega_i^*:
 \quad\text{all eigenvalues are real.}
\]
For $i\ne j$, the relation becomes
\[
 (\omega_i-\omega_j)\braket{\omega_j}{\omega_i}=0.
\]
If $\omega_i\ne\omega_j$, the eigenvectors are orthogonal.
If $\omega_i=\omega_j$, this equation does not impose orthogonality.

\smallheading{Choosing an orthonormal basis within a degeneracy}
Suppose $\omega$ is $m$-fold degenerate. Begin with $m$ linearly
independent eigenvectors in that eigenspace. Gram--Schmidt
orthonormalization replaces them by linear combinations that are
orthonormal. Since these combinations remain in the same eigenspace,
they are still eigenvectors with eigenvalue $\omega$:
\[
 \underbrace{\text{$m$ independent eigenvectors}}_{\text{same eigenvalue}}
 \xrightarrow{\text{Gram--Schmidt}}
 \underbrace{\ket{\omega,1},\ldots,\ket{\omega,m}}_{\text{$m$ orthonormal eigenvectors}}.
\]
The second label distinguishes states within the degeneracy. For
example, a threefold degeneracy has states
$\ket{\omega,1},\ket{\omega,2},\ket{\omega,3}$.

\smallheading{Matrix in an eigenvector basis}
With an orthonormal eigenbasis,
\begin{align*}
 \mel{\omega',k'}{\Omega}{\omega,k}
 &=\omega\braket{\omega',k'}{\omega,k}\\
 &=\omega\,\delta_{\omega'\omega}\delta_{k'k}.
\end{align*}
Different eigenvalues enforce the first delta; orthonormalization
within each eigenspace supplies the second. The matrix is therefore
diagonal:
\[
 [\Omega]_{\mathrm{eigenbasis}}=
 \begin{bmatrix}
 \omega_1&0&\cdots&0\\
 0&\omega_2&\cdots&0\\
 \vdots&\vdots&\ddots&\vdots\\
 0&0&\cdots&\omega_n
 \end{bmatrix}.
\]
For this reason, solving the eigenvalue problem is often called
\emph{diagonalization}.
